% A.45 is reserved by the root for the finite parameter certificate.
\subsection{An exact finite certificate through parameter 31}

\begin{statement}{A.45}{Proved}{Every integral realization from 2 through 31 is recognized}
Every noncuspidal Selmer section in A.14 whose labeled integral
realization has $\lambda_0\in\{2,3,\ldots,31\}$ is a rational
point section.
\end{statement}
\begin{proof}
A.36--A.39 and A.41 leave only
\[
                  12,13,15,16,17,20,21,22,28,29.
\tag{A.45.1}
\]
We supply finite exact certificates for these ten parameters.  All
have potentially good reduction at 2, so their sections are ordinary
there by A.37.  The support and degeneration-label comparison
(A.41.2) remains valid without the odd-valuation hypothesis of
A.41.  In particular, all odd local cusp labels are 0 or 1 and
agree with the prime divisors of $\lambda_0$ or $\lambda_0-1$.

We record two further consequences of the comparisons already
proved.  First, for an odd bad prime $r$, let
$n_r=v_r(\lambda_0)>0$ or $v_r(1-\lambda_0)>0$, as appropriate,
and let $c_r\in\widehat\Z$ be the valuation of the corresponding
completed coordinate.  Then
\[
                     c_r\in n_r\widehat\Z^\times.
\tag{A.45.2}
\]
At an ordinary point, the positive-integer argument in (A.34.2)
gives the stronger equality $c_r=n_r$.  At a cusp, use its
saturated Tate line from A.41.  In Tate bases an integral
intertwiner is upper triangular with unit diagonal entries
$a_\ell,d_\ell$.  Comparing the valuations of the extension
classes gives $2c_{r,\ell}=(a_\ell/d_\ell)2n_r$, up to reversing
the two curves.  Division in $\Q_\ell$ proves (A.45.2) for every
$\ell$, including $\ell=r$.  The splitting extensions here are
unramified quadratic extensions, and the constant $2^8$ in the
cusp parameter is an $r$-adic unit, so no valuation factor is lost.

Second, the exact three-element Galois-set comparison A.39 implies
one global sign $\epsilon\in\{1,-1\}$ with
\[
       \kappa_3(4AB)=
       \epsilon\,\kappa_3(4\lambda_0(1-\lambda_0))
                 \quad\text{in }H^1(\Q,\mu_3).
\tag{A.45.3}
\]
To check the sign, restrict to $K=\Q(\zeta_3)$.  Both sets of
cube roots are torsors for the constant group $\F_3$.
Every permutation of a three-element set is an affine map
$x\mapsto\epsilon x+b$, so an equivariant bijection identifies
their translation cocycles up to that one sign.  Restriction on
Kummer classes is injective: if a rational number $d$ becomes a
cube in $K$, its norm makes $d^2$ a rational cube, whence $d$ is
itself a rational cube.  This proves (A.45.3) over $\Q$, rather
than allowing an independently chosen sign at each prime.

Here is the complete small calculation at 2.  Put
$a=v_2(x_2)$, $b=v_2(1-x_2)$ and let $u_A,u_B\in\{1,3,5,7\}$
be the unit parts modulo 8 of the squareclasses of
$\lambda_0,1-\lambda_0$.  The possibilities are
$(a,b)=(m,0),(0,m),(-m,-m)$.  The Legendre formula gives
$v_2(j(x_2))=8-2m$, so potential good reduction gives $1\leq m\leq4$.
The valuation parities must match the two fixed squareclasses.
In the three cases the unit conditions are respectively
\[
 u_B\equiv1-2^m u_A,\qquad
 u_A\equiv1-2^m u_B,\qquad
 u_B\equiv2^m-u_A\pmod8.
\tag{A.45.4}
\]
Exhausting these four integers gives
\[
\begin{array}{c|c|l}
 \lambda_0&(u_A,u_B)&(a,b)\\ \hline
 12&(3,5)&(2,0),\ (-4,-4)^*\\
 13&(5,5)&(0,2),\ (2,0)\\
 15&(7,1)&(0,1)\\
 20&(5,5)&(0,2),\ (2,0)\\
 21&(5,3)&(0,2),\ (-4,-4)^*\\
 22&(3,3)&(1,0)\\
 28&(7,5)&(2,0),\ (-2,-2)\\
 29&(5,1)&(0,2),\ (4,0)^*.
\end{array}
\tag{A.45.5}
\]
The starred alternatives are excluded by (A.45.3).  For parameters
$12,21,29$, the valuation of $4\lambda_0(1-\lambda_0)$ at 2 is
$4\equiv1\pmod3$, whereas the starred alternatives give
$2+a+b\equiv0\pmod3$.  They cannot agree with either sign.
The table and every logarithm valuation used below are recomputed
using integer arithmetic in \texttt{code/legendre\_support\_checks.py}.

For $12$ and $29$ the surviving row gives a coordinate with zero
2-valuation supported at the single odd prime 11 or 29.  For
$13$ and $20$, either surviving pair gives such a coordinate,
supported respectively at one of $3,13$ or one of $5,19$.
The one-prime-support argument (A.41.3)--(A.41.6) recognizes each
of these four parameters: a section with both cusp labels would
force a primary exponent to vanish, contradicting multiplicative
inertia at its supporting prime.

For $16$, the class $A$ is supported only at 2; for $17$, the
class $B$ is supported only at 2.  Their local sections at 2 are
ordinary, so the respective valuation is an honest integer.
The class is therefore an actual rational number.  If $A=a_0$,
local Kummer injectivity places every ordinary centre at $a_0$;
if $B=b_0$, it places every ordinary centre at $1-b_0$.
Including the three boundary points gives finite support in either
case.  Finite-support recognition makes the noncuspidal section
a rational point section.

For $15$, (A.45.5) and the support calculation give
\[
                    A=3^\alpha5^\gamma,
                    \qquad B=-2\,7^\beta.
\]
A cusp 0 at 5 would require
$\log_5(2)+\beta_5\log_5(7)=0$.
But
\[
 v_5(\log_5 2)=v_5(2^4-1)=1,\qquad
 v_5(\log_5 7)=v_5(7^4-1)=2,
\]
so it would require $\beta_5$ to have valuation $-1$, impossible
for $\beta_5\in\Z_5$.  Thus the section at 5 is ordinary, and
(A.45.2) in its ordinary form gives $\gamma=1$.
Valuation at 5 in (A.45.3) forces $\epsilon=1$; valuation at 7
then gives $\beta\equiv1\pmod3$.  A cusp 0 at 3 would require
$-2\,7^\beta$ to be a cube in $\Q_3$.  The cubes of units modulo
9 are $1,8$, and $7$ and $2$ have the same class modulo cubes.
This condition instead gives $1+\beta\equiv0\pmod3$.
Thus no cusp 0 occurs at 3 or 5.  The only remaining cusp label is
1, and A.37 recognizes the section.

For $21$, the surviving 2-adic contact pair gives
\[
                         A=3^\alpha7^\gamma.
\]
Both $\alpha_5,\gamma_5$ are units by (A.45.2).
A cusp 1 at 5 would give
$\alpha_5\log_5(3)+\gamma_5\log_5(7)=0$.
The two summands have valuations one and two, since
$3^4-1=80$ and $7^4-1=2400$.  Their sum cannot vanish.
No cusp 1 remains, so A.37 applies.

For $22$, the contact pair $(1,0)$ gives
$B=-3^\beta7^\gamma$, with $\beta_{11},\gamma_{11}$ units by
(A.45.2).  A cusp 0 at 11 would give
$\beta_{11}\log_{11}(3)+\gamma_{11}\log_{11}(7)=0$.
Here
\[
 v_{11}(3^{10}-1)=2,
             \qquad v_{11}(7^{10}-1)=1.
\]
Again the unique smaller valuation prevents cancellation.
This excludes the only cusp-0 prime and A.37 applies.

Finally consider $28$.  For contact pair $(2,0)$ the class
$B=-3^\beta$ has one-prime support, so A.41's argument applies.
For contact pair $(-2,-2)$, the exact classes are
\[
                 A=2^{-2}7^\alpha,
                 \qquad B=-2^{-2}3^\beta.
\]
A section not already recognized would have cusp 1 at 3 and cusp 0
at 7.  The first forces $7^\alpha=4$ in the local completed group
at 3.  Modulo cubes this gives $\alpha\equiv2\pmod3$, since
$7\equiv-2$ modulo 9 and $4=2^2$.  On the other hand,
$v_2(4AB)=-2\equiv1\pmod3$ and
$v_2(4\cdot28\cdot(-27))=4\equiv1\pmod3$, so (A.45.3)
forces $\epsilon=1$.  Its valuation at 7 gives
$\alpha\equiv1\pmod3$, a contradiction.
This completes every case in (A.45.1) and proves the statement.
\end{proof}
